% ===================================================================
%  HOTO Na 6 — Cikin gida, kayan gida, abinci, haske.
%
%  Traduction, faite en 2026, en haoussa d'aujourd'hui, sur l'ido de
%  texto/io. Regles de la colonne : voir 10-hoto-01.tex. Cinq scenes,
%  cinq numerotations. C'est le tableau le plus charge du livret :
%  211 renvois.
%
%  LA CUISINE EST HAOUSSA ET LE SALON NE L'EST PAS, ET LES DEUX SONT
%  SUR LA MEME PLANCHE. C'est le partage le plus net du tableau, et il
%  se lit d'une piece a l'autre sans une exception :
%
%    murhu (1)      le foyer        tukunya (43)   la marmite
%    ludayi (21)    la louche       tanda (20)     le four
%    miya (24)      la sauce        gishiri (14)   le sel
%    barkono (15)   le poivre       gawayi (9)     le charbon
%    tsintsiya (34) le balai        mazurari (38)  l'entonnoir
%    matata (28)    la passoire     ashana (12)    les allumettes
%
%  Contre, au salon : sofa, firam, albam, fiyano, darduma, lambarkin,
%  fitilar lantarki. Pas un objet de cette piece-la n'existait ici
%  avant qu'on l'y porte, et la langue le dit en une syllabe. La
%  cuisine, elle, n'a rien emprunte du tout : on y fait bouillir
%  depuis toujours, et le mot du feu n'a pas change.
%
%  ET LE CHARBON DE TERRE (17) EST LA MEILLEURE PREUVE : il se dit
%  « gawayin dutse », le charbon de bois de pierre. Le haoussa avait
%  le gawayi — le charbon de bois, qu'on fait au village — et il a
%  nomme l'autre par la matiere. Ni emprunt ni calque : une
%  DERIVATION, sur un mot qu'on avait deja.
%
%  CINQUIEME REFUS DE LA COLONNE : LE POT A EAU (3). Le haoussa a
%  « randa », et c'est le mot qui vient d'abord : la randa est la
%  grande jarre de terre poreuse qui se tient dans l'entree de toute
%  maison haoussa et qui garde l'eau fraiche par evaporation. La
%  planche grave un broc de faience pose sur une table de toilette
%  europeenne, a cote d'une cuvette. On ecrit « tulun ruwa ». Meme
%  raison pour la coupe a fruits (10), ou « kwarya », la calebasse,
%  aurait ete le mot le plus naturel du monde : la planche montre un
%  compotier a pied. Une matiere n'est pas une forme.
%
%  DEUX LAMPES DANS LE MEME TABLEAU, ET UN SEUL MOT : « fitila » est
%  la lampe a petrole du salon (12) — « fitilar kananzir » — et la
%  lampe electrique (39) — « fitilar lantarki » —, comme c'etait deja
%  la lampe a huile de l'atelier de dessin au tableau 1. Le haoussa ne
%  renomme pas la lumiere quand le combustible change ; il ajoute
%  seulement de quoi elle brule. Trois siecles de progres tiennent
%  dans un complement de nom.
%
%  UNE QUATRIEME PAIRE D'HOMONYMES, APRES kaka, kwano ET allo :
%  « allura » est l'aiguille a coudre (37) de la corbeille a ouvrage
%  et c'etait deja l'aiguille de l'horloge au tableau 3 — « ƙaramar
%  allura », la petite aiguille, sur le chiffre X. Le haoussa nomme
%  par la FORME et ne se soucie pas de l'usage : ce qui est long,
%  mince et pointu est une allura, qu'il perce le tissu ou qu'il
%  marque l'heure.
%
%  UNE FAUTE D'ORDRE, LA DEUXIEME DE LA COLONNE, ET C'EST LA MEME
%  ESPECE QU'AU TABLEAU 3. Au bloc t06-c3-02-1 la fontaine (3) avait
%  ete posee avant l'escabeau (4), alors que l'ido dit « pendas, super
%  la skabelo, fonteneto orniva » : le lieu d'abord, l'objet ensuite.
%  Exactement la phrase des roseaux et du feu d'artifice au tableau 3,
%  exactement le meme glissement. DEUX FOIS LA MEME CAUSE : le haoussa
%  est a tete initiale et pose son complement de lieu apres, et rien
%  ne l'y oblige — « a saman ƙaramar kujera » s'insere tres bien
%  devant l'objet. La colonne a un seul defaut et il est constant.
%
%  ET LE DE A COUDRE EST « ƙoƙon yatsa », LA CALEBASSE DU DOIGT. Le
%  meme « ƙoƙo » — la coque, la demi-calebasse — nommait le crane au
%  tableau 2, « ƙoƙon kai ». Un mot de coque pour tout ce qui coiffe
%  et protege, et deux objets qui n'ont rien a voir sauf leur forme.
% ===================================================================

%%K t06-apar-1 apar
\VUsaut{6.0mm}
\VUtitre{15.0pt}{60}{HOTO Na 6}
\VUsaut{1.4mm}
\VUtitre{11.4pt}{0}{Cikin gida, kayan gida, abinci,}
\VUsaut{0.4mm}
\VUtitre{11.4pt}{0}{haske.}
\VUsaut{2.2mm}

%%K t06-apar-2 apar
Gidan ya cika. Baffan Ioannes ya zauna a sabon gidansa, wanda muka san
tsarinsa. Yanzu mu ga yadda ya shirya kayan gidansa. Da farko za mu
bincika :

%%K t06-tit-1 sub
\VUpk{10.2pt}{40}{Ɗakin kwana.}
\VUsaut{3.0mm}

%%K t06-c1-01-1 p
1. --- Mun iso a lokacin da bai dace ba, domin mai kula da ɗaki tana
cikin tsaftace \VUgras{ɗakin.}

%%K t06-c1-02-1 p
2. --- A kan \VUgras{teburin ado} \textsuperscript{(1)} ga kayayyaki iri
iri nan da can, waɗanda ake wanke kai da su, ake yin kitso da su, a
taƙaice ake yin ado da su. Su ne \VUgras{kwanon wanka}
\textsuperscript{(2)} da \VUgras{tulun ruwa} \textsuperscript{(3)}, da
\VUgras{buroshin kai} \textsuperscript{(4)}, da \VUgras{matsefi}
\textsuperscript{(5)}, da \VUgras{sabulu} \textsuperscript{(6)} da
\VUgras{soso} \textsuperscript{(7)}, a ƙarshe kuma \VUgras{ƙaramar
kwalaba} \textsuperscript{(8)} ta maganin haƙori na ruwa.

%%K t06-c1-03-1 p
3. --- A ƙasa, gaban teburin ado, ga \VUgras{guga}
\textsuperscript{(9)} domin a tara ruwa mai ƙazanta.

%%K t06-c1-04-1 p
4. --- Muna gani a \VUgras{ɗakin gaba} \textsuperscript{(10)} waɗansu
\VUgras{ƙugiyoyin tufafi} \textsuperscript{(13)} da \VUgras{laima}
\textsuperscript{(11)} biyu cikin \VUgras{mariƙin laima}
\textsuperscript{(12)}.

%%K t06-c1-05-1 p
5. --- A ɗaya gefen ƙofar akwai \VUgras{kabad mai madubi}
\textsuperscript{(14)} wanda yake ɗauke da \VUgras{farar kaya}
\textsuperscript{(15)}.

%%K t06-c1-06-1 p
6. --- Mu yi amfani da lokacin da \VUgras{mai kula da ɗaki}
\textsuperscript{(6)} take shirya \VUgras{gado} \textsuperscript{(17)},
domin mu bincika sassansa iri iri. \VUgras{Katakon gado}
\textsuperscript{(20)} yana ɗauke da \VUgras{katifar ƙarfe}
\textsuperscript{(21)} mai kyau wadda Iustina za ta ajiye a kanta
\VUgras{katifar ulu} \textsuperscript{(22)} nan take. Sa’an nan za ta
ɗauko daga kujeru \VUgras{zannuwan gado} \textsuperscript{(23)} biyu da
\VUgras{bargo} \textsuperscript{(24)}. Sa’an nan za ta ajiye a kan kan
gadon \VUgras{dogon matashi} \textsuperscript{(25)} da \VUgras{matashin
kai} \textsuperscript{(26)} waɗanda suke tare da \VUgras{bargon gashi}
\textsuperscript{(27)} a kan \VUgras{tabarmar gado}
\textsuperscript{(32)}. Tana amfani da gungun gashin tsuntsu domin ta
kakkaɓe ƙura daga \VUgras{labule} \textsuperscript{(19)} da \VUgras{rufin
gado} \textsuperscript{(18)}, ga wani ɓangaren aikin ya ƙare kenan.

%%K t06-c1-07-1 p
7. --- Sa’an nan za ta yi gyara kaɗan ga \VUgras{ƙaramin teburin
gado} \textsuperscript{(28)}, ta murza \VUgras{agogon farkarwa}
\textsuperscript{(29)}, ta shirya \VUgras{ƙaramar fitilar dare}
\textsuperscript{(30)}, ta ɗauke \VUgras{kyandir} \textsuperscript{(31)}
domin ta tsaftace mariƙinsa.

%%K t06-tit-2 sub
\VUsaut{5.0mm}
\VUpk{10.2pt}{40}{Kwandon aiki da kayan murhu.}
\VUsaut{3.0mm}

%%K t06-c1-08-1 p
8. --- \VUgras{Kwandon aikin hannu} \textsuperscript{(34)} kusa da
\VUgras{ƙaramar kujera} \textsuperscript{(33)} shi ma yana buƙatar gyara
ƙwarai. Za ta naɗe \VUgras{ɗinkin ado} \textsuperscript{(35)}, ta sa
cikin \VUgras{teburin ɗinki} \textsuperscript{(38)} \VUgras{ƙoƙon
yatsa,} da \VUgras{zare} \textsuperscript{(36)}, da \VUgras{allura}
\textsuperscript{(37)} da \VUgras{almakashi} \textsuperscript{(39)}, ta
kuma rufe murfin teburin da \VUgras{mabuɗi} \textsuperscript{(40)}.

%%K t06-c1-09-1 p
9. --- A kan \VUgras{teburin ƙafa ɗaya} \textsuperscript{(41)} na
tsakiya, Iustina za ta sabunta ruwan \VUgras{kwalaban ruwa}
\textsuperscript{(42)}. Za ta sa \VUgras{sukari} cikin \VUgras{akwatin
sukari} \textsuperscript{(43)}, ta tsaftace \VUgras{matsin sukari}
\textsuperscript{(44)} ta kuma ɗauke \VUgras{buroshin tufafi}
\textsuperscript{(46)}.

%%K t06-c1-10-1 p
10. --- Gefen dama na ɗakin ba zai buƙaci aiki mai yawa ba a gare
ta, domin \VUgras{doguwar kujera} \textsuperscript{(47)} da
\VUgras{matashinta} \textsuperscript{(48)} suna a wurarensu daidai,
kuma, tun da ba sanyi ba ne, ba a motsa komai a cikin kayan
\VUgras{murhun bango} \textsuperscript{(49)}. \VUgras{Shebur}
\textsuperscript{(50)}, da \VUgras{manyan matsi} \textsuperscript{(51)},
da \VUgras{sandunan wuta} \textsuperscript{(55)}, da \VUgras{shingen
toka} \textsuperscript{(56)} suna da tsabta ; ba a motsa \VUgras{shingen
wuta} \textsuperscript{(54)} ba, \VUgras{akwatin itace}
\textsuperscript{(52)} kuma yana cike da \VUgras{icen wuta}
\textsuperscript{(53)}.

%%K t06-noto-1 noto
\VUnotes{143.2mm}{%
(*) Babo = jariri, E. baby, F. bébé, I. bambino.%
}

%%K t06-noto-2 noto
\VUnotes{143.2mm}{%
(*) Lambrequino = F. lambrequin, S. lambrequines, Port. lambrequins,
har ma D. E. lambrequins, don haka D. E. F. P. S.%
}

%%K t06-c1-11-1 p
11. --- Duk da haka za ta kakkaɓe ƙura daga \VUgras{agogo mai lilo}
\textsuperscript{(57)}, da \VUgras{mariƙan kyandir}
\textsuperscript{(58)} da \VUgras{kwanonin ado} \textsuperscript{(59)},
a ƙarshe kuma ta share \VUgras{madubi} \textsuperscript{(60)}.

%%K t06-c1-12-1 p
12. --- Game da \VUgras{shimfiɗar} \textsuperscript{(61)} jaririn (jariri
(*)), ta riga ta yi shiri.

%%K t06-tit-3 sub
\VUsaut{5.0mm}
\VUpk{10.2pt}{40}{Falo.}
\VUsaut{3.0mm}

%%K t06-c2-01-1 p
1. --- Yanzu ga falo. \VUgras{Ɗaki} ne mai kyau ƙwarai. Murhun bango,
wanda yake a kusurwar dama, an yi masa ado da \VUgras{ƙaramin
mutum-mutumi} \textsuperscript{(1)} mai laushi, da \VUgras{tuluna}
\textsuperscript{(3)} biyu masu kyau, an kuma lulluɓe shi da
\VUgras{lambarkin} \textsuperscript{(2)} na kaɗi (*).

%%K t06-c2-02-1 p
2. --- Domin a hana tartsatsin wuta ƙona darduma, an ajiye
\VUgras{shingen tartsatsi} \textsuperscript{(4)} na jan ƙarfe a gaban
murhun.

%%K t06-c2-03-1 p
3. --- A kusa, \VUgras{teburin rubutu} \textsuperscript{(5)} na mata
mai kwarjini yana a buɗe. A kansa \VUgras{uwar gida}
\textsuperscript{(23)} take rubuta wasiƙunta.

%%K t06-c2-04-1 p
4. --- An yi wa taga ado da \VUgras{labulen taga}
\textsuperscript{(6)} tare da \VUgras{ɗauransu} \textsuperscript{(8)} da
aka rataya a \VUgras{ƙugiyoyi} \textsuperscript{(7)} da kuma manyan
labule na kaɗi \VUgras{waɗanda aka naɗe a sama}
\textsuperscript{(9)}.

%%K t06-c2-05-1 p
5. --- A cikin ramin taga, \VUgras{mariƙin tukunya}
\textsuperscript{(11)} yana ɗauke da koren \VUgras{tsiro}
\textsuperscript{(10)}, itacen dabino mai kyau wanda ganyensa suke
miƙewa kusan har \VUgras{fitilar kananzir} \textsuperscript{(12)}.

%%K t06-c2-06-1 p
6. --- \VUgras{Murfin fitilar} \textsuperscript{(13)} Maria ta shirya
shi, \VUgras{budurwa} \textsuperscript{(14)} mai kwarjini wadda take kan
fiyano \textsuperscript{(15)}. A zaune sosai a kan \VUgras{ƙaramar
kujera} \textsuperscript{(16)}, tana nazarin waƙa. Tana da gwaninta da
kula ƙwarai, kamar yadda \VUgras{kabad na waƙarta}
\textsuperscript{(17)} yake shaida, wanda aka shirya sarai.

%%K t06-tit-4 sub
\VUsaut{5.0mm}
\VUpk{10.2pt}{40}{Ɗan iska.}
\VUsaut{3.0mm}

%%K t06-c2-07-1 p
7. --- Ɗan uwanta, Paulus, yana neman \VUgras{littafi}
\textsuperscript{(19)} a cikin \VUgras{kabad na littattafai}
\textsuperscript{(18)}. \VUgras{Saurayi} \textsuperscript{(20)} ne, mai
motsi kuma marar haƙuri. Yana rataya kansa a kabad ɗin domin ya kai ga
littafin da yake can sama, yana cikin haɗarin faɗar da
\VUgras{mutum-mutumi} \textsuperscript{(21)} da yake saman.

%%K t06-c2-08-1 p
8. --- Yana amfani da lokacin da mahaifiyarsa take hira da wata
abokiya, \VUgras{mata} \textsuperscript{(22)} wadda ta zo ta ɗan yi
kwana kaɗan a gidanta, domin ya yi wannan wauta. Mahaifiyar tana
zaune a kan \VUgras{sofa} \textsuperscript{(24)} kusa da \VUgras{babbar
kujera} \textsuperscript{(25)}, abokiyarta kuwa tana kan
\VUgras{kujera} \textsuperscript{(27)} wadda ba a ganin komai daga
gare ta sai \VUgras{madogarar bayanta} \textsuperscript{(26)}.

%%K t06-c2-09-1 p
9. --- Babban \VUgras{firam} \textsuperscript{(30)} da yake rataye a
bango yana ɗauke da \VUgras{hoton} \textsuperscript{(29)} wata ƴar uwa.
A cikin \VUgras{littafin hotuna} \textsuperscript{(32)} da aka ajiye a
kan tebur an tara \VUgras{hotunan} \textsuperscript{(33)} sauran ƴan
uwan gidan. Yaran sun gama sassaƙe shafukansa yanzu, domin
\VUgras{kujeru} \textsuperscript{(31)} har yanzu suna tare kewaye da
teburin.

%%K t06-c2-10-1 p
10. --- \VUgras{Tulun ƙasar Sin} \textsuperscript{(34)} na fenti,
wanda yake kusa da \VUgras{wuƙar takarda} \textsuperscript{(35)}, an ba
Maria shi a ranar bikinta, \VUgras{shingen ado}
\textsuperscript{(36)} kuwa wanda yake yi wa kusurwar ɗakin ado, baffanta
ne ya kawo shi daga ƙasar Sin.

%%K t06-c2-11-1 p
11. --- Da \VUgras{zanuka} \textsuperscript{(37)} masu kyau ƙwarai
waɗanda aka rataya a \VUgras{bango} \textsuperscript{(42)} suke faranta
shi, da \VUgras{fitilar rufi} \textsuperscript{(38)} take haskaka shi,
wadda \VUgras{fitilun lantarkinta} \textsuperscript{(39)} suke haske da
farin ciki da yamma, wannan falo yana da kyau ƙwarai. Mai laushin
\VUgras{darduma} \textsuperscript{(40)} da aka shimfiɗa a kan katakon
ƙasa, da \VUgras{ƙararrawar lantarki} \textsuperscript{(41)} mai saukin
amfani, sun cika jin daɗinsa.

%%K t06-tit-5 sub
\VUsaut{3.6mm}
\VUpk{10.2pt}{40}{Ɗakin cin abinci.}
\VUsaut{3.0mm}

%%K t06-c3-01-1 p
1. --- Ƙararrawar abincin dare ta kaɗa. Dukan iyalin, sai Maria, sun
taru kewaye da tebur. An ɗumama ɗakin cin abinci da kyau da
\VUgras{murhu} \textsuperscript{(1)} na yumɓu mai kyau ƙwarai, mai
sirarin \VUgras{bututu} \textsuperscript{(2)}.

%%K t06-c3-02-1 p
2. --- Wannan ɗaki ba shi da kayan gida masu yawa. A gefen bango, a
saman \VUgras{ƙaramar kujera} \textsuperscript{(4)}, akwai
\VUgras{ƙaramar marmaron ado} \textsuperscript{(3)} a rataye. Kusa ƙwarai, akwai
\VUgras{kabad na kwanuka} \textsuperscript{(5)} wanda ake ajiye a kansa
abincin sanyi, da kwanonin kayan zaki, da kayan tebur iri iri waɗanda
ba a buƙatarsu nan take.

%%K t06-c3-03-1 p
3. --- Haka nan muna gani a kan kantar sama, \VUgras{gasashshen ɗan
kaza} \textsuperscript{(7)}, da \VUgras{kifi} \textsuperscript{(8)} da
\VUgras{babban kwano} \textsuperscript{(10)} na \VUgras{ƴaƴan itace}
\textsuperscript{(9)}. A ƙasa ga \VUgras{cuku} \textsuperscript{(11)}, da
\VUgras{burodi} \textsuperscript{(12)}, da \VUgras{mariƙin kwalabe}
\textsuperscript{(13)} wanda yake ɗauke da duk abin da ake buƙata domin
a yi wa salatin kayan ƙamshi : mai, da ruwan tsami, da
\VUgras{gishiri} \textsuperscript{(14)} da \VUgras{barkono}
\textsuperscript{(15)}. A ƙarshe, a kan kantar ƙasa akwai
\VUgras{kek} \textsuperscript{(17)} kusa da \VUgras{mariƙin mastad}
\textsuperscript{(16)}.

%%K t06-c3-04-1 p
4. --- Agogon ɗakin cin abinci, wanda ake kira \VUgras{agogon bango}
\textsuperscript{(20)}, yana da siffa ta musamman ƙwarai. An sa shi
cikin firam na katako wanda yake kuma ɗauke da \VUgras{ma’aunin
yanayi} \textsuperscript{(19)} da \VUgras{ma’aunin zafi}
\textsuperscript{(18)}.

%%K t06-tit-6 sub
\VUsaut{4.6mm}
\VUpk{10.2pt}{40}{Hidimar abincin dare.}
\VUsaut{3.0mm}

%%K t06-c3-05-1 p
5. --- A dukan bangayen ɗakin an liƙa \VUgras{katakon bango}
\textsuperscript{(21)} mai ƙyalli. \VUgras{Labulen ƙofa}
\textsuperscript{(22)} na kaɗi yana rufe ƙofar da take kai wa baranda
sosai. A can ne iyalin za su je su sha kofi bayan abincin dare. Tuni
\VUgras{kofuna} \textsuperscript{(24)} da \VUgras{farantan kofuna}
\textsuperscript{(25)} suke a shirye a kan \VUgras{teburin kwanuka}
\textsuperscript{(23)}.

%%K t06-c3-06-1 p
6. --- A saman tebur an kafa a silin \VUgras{fitila mai rataya}
\textsuperscript{(26)} wadda fitilarta mai haske ƙwarai take haskaka
dukan ɗakin cin abinci da yamma.

%%K t06-c3-07-1 p
7. --- Yanzu mu bincika \VUgras{teburin cin abinci}
\textsuperscript{(27)} da kansa. Lokacin da iyalin suka shigo, kayan
cin abinci sun riga sun shirya. A kan \VUgras{zanen tebur}
\textsuperscript{(28)} an ajiye, a wurin kowane mai ci,
\VUgras{faranti} \textsuperscript{(33)}, da \VUgras{ɗan tawul}
\textsuperscript{(29)}, da \VUgras{cokali} \textsuperscript{(34)}, da
\VUgras{cokali mai yatsa} \textsuperscript{(35)}, da \VUgras{wuƙa}
\textsuperscript{(36)} da \VUgras{gilashi} \textsuperscript{(32)}. Akwai
kuma \VUgras{kwalabe} \textsuperscript{(30)} na giya da \VUgras{kwalaban
ruwa} \textsuperscript{(31)}.

%%K t06-c3-08-1 p
8. --- \VUgras{Mai hidima} \textsuperscript{(39)} ya fara kawo
\VUgras{kwanon miya} \textsuperscript{(37)}, wanda miya kaɗan take har
yanzu tana hayaƙi a ciki. Yanzu yana ba da \VUgras{kwano}
\textsuperscript{(38)} na farko, wato kwanon nama. Sa’an nan zai
gabatar da waɗanda muka riga muka ambata, a ƙarshe kuma, bayan
\VUgras{kwanon kayan zaki,} zai yi amfani da lokacin da iyayen gidansa
suka shiga baranda, domin ya ɗauke kayan tebur ya kuma sake shirya
ɗakin cin abinci.

%%K t06-c3-08-2 p
\textit{Bayani.} --- \textit{An nuna kalmomin da ake amfani da su
game da abinci a nan a taƙaice ƙwarai. Za a cika jerin a kan hotuna
biyu, « Otal » (na 13) da « Kasuwa » (na 15).}

%%K t06-tit-7 sub
\VUsaut{4.6mm}
\VUpk{10.2pt}{40}{Ɗakin girki.}
\VUsaut{3.0mm}

%%K t06-c4-01-1 p
1. --- A ɗakin girki, wanda muke kallo ta ƙofar da take rabin buɗe,
muna ganin rikici mai ban sha’awa. A gaban \VUgras{murhu}
\textsuperscript{(1)} inda \VUgras{wuta} \textsuperscript{(3)} take ci da
ƙarfi, an kafa \VUgras{sandar juya nama} \textsuperscript{(5)}.
\VUgras{Gasashshen nama} \textsuperscript{(7)} mai kyau \VUgras{an tsira
shi da sanda} \textsuperscript{(6)} yana dahuwa.

%%K t06-c4-02-1 p
2. --- \VUgras{Abin busa wuta} \textsuperscript{(8)} da \VUgras{shebur
na gawayi} \textsuperscript{(9)}, waɗanda aka yi amfani da su yanzu, sun
zauna a ƙasa kamar yadda \VUgras{icen wuta} \textsuperscript{(10)} suka
zauna.

%%K t06-c4-03-1 p
3. --- An shirya saman murhun ; \VUgras{fitilar hannu}
\textsuperscript{(11)}, da \VUgras{akwatin ashana}
\textsuperscript{(13)}, inda \VUgras{ashana} \textsuperscript{(12)} suke,
da \VUgras{mariƙin kyandir} \textsuperscript{(14)} da \VUgras{ƙaramin
mariƙin kyandir} \textsuperscript{(15)} tare da \VUgras{kyandir}
\textsuperscript{(16)} suna a wurarensu. Amma babbar \VUgras{gugar
gawayi} \textsuperscript{(18)}, cike da \VUgras{gawayin dutse}
\textsuperscript{(17)} wanda \VUgras{mai girki} \textsuperscript{(23)} ta
cika yanzu daga ɗakin ƙasa, ta zauna a tsakiyar ɗakin girkin.

%%K t06-c4-04-1 p
4. --- Hakika, matar nan mai tausayi dole ta yi gaggawa domin abincin
rana ya shirya a kan lokaci. Yanzu tana kusa da \VUgras{murhun girki}
\textsuperscript{(19)} tana juya \VUgras{miya} \textsuperscript{(24)}
wadda bai kamata ta yi ƙona ba.

%%K t06-c4-05-1 p
5. --- A cikin \VUgras{tanda} \textsuperscript{(20)} kek yana toyuwa,
wanda ta shirya domin ƙaramin abincin yaran. A saman murhun girki
akwai \VUgras{ludayi} \textsuperscript{(21)} da \VUgras{ludayi mai rami}
\textsuperscript{(22)} a rataye.

%%K t06-tit-8 sub
\VUsaut{4.0mm}
\VUpk{10.2pt}{40}{Kwanuka.}
\VUsaut{3.0mm}

%%K t06-c4-06-1 p
6. --- \VUgras{Kayan girki} \textsuperscript{(25)}, waɗanda aka rataya
a ƙarshen ɗakin, suna da tsabta ƙwarai. Kyawawan \VUgras{tukwane na
jan ƙarfe} \textsuperscript{(26)}, da \VUgras{tulun madara}
\textsuperscript{(27)}, da \VUgras{matata} \textsuperscript{(28)} da
\VUgras{abin nika cuku} \textsuperscript{(29)} an tsaftace su da kula.

%%K t06-c4-07-1 p
7. --- An ajiye \VUgras{akwatin gishiri} \textsuperscript{(30)} a kan
kabad na kwanuka kusa da \VUgras{tukunyar shayi} \textsuperscript{(31)}
da \VUgras{babbar kwalaban mai} \textsuperscript{(32)}. Wataƙila
\VUgras{aljihun tebur} \textsuperscript{(33)} yana ɗauke da kayan ci da
wuƙaƙen ɗakin girki.

%%K t06-c4-08-1 p
8. --- \VUgras{Tsintsiya} \textsuperscript{(34)} da \VUgras{gungun
gashin tsuntsu} \textsuperscript{(35)} suna a kusurwa. Mai girki ta yi
amfani da \VUgras{gungumen katako} \textsuperscript{(36)} yanzu, ta bar
shi kusa da taga.

%%K t06-c4-09-1 p
9. --- \VUgras{Tawul ɗin hannu} \textsuperscript{(37)} yana rataye a
bango, kusa da \VUgras{mazurari} \textsuperscript{(38)}. A wannan
ɓangaren ɗakin girkin an ajiye \VUgras{gasa} \textsuperscript{(39)}, da
\VUgras{wuƙar yankan nama} \textsuperscript{(40)} da \VUgras{katakon
yankan nama} \textsuperscript{(41)}. Sa’an nan, a saman wuƙar, a kan
\VUgras{kanta} \textsuperscript{(42)}, akwai \VUgras{tukwane}
\textsuperscript{(43)}, da \VUgras{tukunyar tafasa}
\textsuperscript{(44)} tare da \VUgras{murfinta} \textsuperscript{(45)}
da \VUgras{babbar tukunya} \textsuperscript{(46)}. \VUgras{Kaskon soya}
\textsuperscript{(47)} yana rataye a kusa.

%%K t06-c4-10-1 p
10. --- \VUgras{Famfo} \textsuperscript{(48)} na jan ƙarfe kaɗai yake
haskakawa a wannan kusurwa. \VUgras{Wurin wanke kwanuka}
\textsuperscript{(49)}, wanda aka ajiye \VUgras{babban tulu}
\textsuperscript{(50)} mai kauri a kansa, wataƙila yana da saukin
amfani ƙwarai da ƙaramin kabad ɗinsa, inda ake ajiye \VUgras{gangar
ruwan tsami} \textsuperscript{(51)} tare da \VUgras{famfonta}
\textsuperscript{(52)}, da \VUgras{akwatin shara} \textsuperscript{(53)}
da \VUgras{faranti masu ƙazanta} \textsuperscript{(54)}.

%%K t06-tit-9 sub
\VUsaut{4.0mm}
\VUpk{10.2pt}{40}{Waɗansu kaya da aka ajiye a ɗakin girki.}
\VUsaut{3.0mm}

%%K t06-c4-11-1 p
11. --- \VUgras{Babban kwano} \textsuperscript{(55)} da ake wanke
\VUgras{ganyayyaki} \textsuperscript{(58)} a ciki da \VUgras{tukunyar
ƙarfe} \textsuperscript{(56)} da aka ajiye a kan \VUgras{ƙasan tayal}
\textsuperscript{(57)} wataƙila suna da wurarensu a wannan kabad ma.

%%K t06-c4-12-1 p
12. --- Lalle mai girki ba ta rasa murhuna, domin ga kuma, a kusurwar
tebur, \VUgras{murhun iskar gas} \textsuperscript{(59)} wanda ruwa yake
zafi a kansa cikin ƙaramar tukunya. \VUgras{Tururi}
\textsuperscript{(60)} da yake fitowa daga cikinta yana nuna mana cewa
wataƙila tana tafasa. Kila za a yi amfani da wannan ruwa domin kofi,
domin ga shi a ɗayan ƙarshen teburin, \VUgras{tulun kofi}
\textsuperscript{(64)} da \VUgras{niƙar kofi} \textsuperscript{(63)}.

%%K t06-c4-13-1 p
13. --- An haɗa murhun da \VUgras{bakin iskar gas}
\textsuperscript{(62)} ta wurin doguwar \VUgras{ƙaramar bututun roba}
\textsuperscript{(61)}.

%%K t06-c4-14-1 p
14. --- \VUgras{Mai aikin yini} \textsuperscript{(66)} wadda take zuwa
ta taimaki mai girki, tana shirya ganyayyaki domin abincin dare. Idan
an tsaftace su aka kuma wanke su, za ta sa su cikin \VUgras{kwano}
\textsuperscript{(65)} tana jiran lokacin dahuwarsu.

%%K t06-tit-10 sub
\VUsaut{4.0mm}
{\centering\fontsize{10.2pt}{10.2pt}\selectfont\textls[40]{\scshape Ɗakin Wanka.} \textit{(Ɗakin yin wanka.)}\par}
\VUsaut{3.0mm}

%%K t06-c5-01-1 p
1. --- Abu ɗaya kaɗai ya rage mana mu yi na ɓoye domin mu san manyan
ɗakunan gidan : mu ga tsarin ɗakin wanka. Wannan ba zai ɗauki lokaci
mai tsawo ba, domin ba babba ba ne, kuma za mu san da sauri cewa
\VUgras{kwanon wanka} \textsuperscript{(1)} yana da \VUgras{abin ɗumama
ruwa} \textsuperscript{(2)} mai kyau, cewa \VUgras{kwanon sabulu}
\textsuperscript{(3)} yana inda hannun mai wanka ya isa kuma cewa
\VUgras{mariƙin tawul} \textsuperscript{(5)} lalle zai sa
\VUgras{tawul} \textsuperscript{(4)} su bushe cikin sauƙi.

%%K t06-c5-02-1 p
2. --- Wannan ɗaki yana da bangayensa a rufe da \VUgras{tayal}
\textsuperscript{(6)}, waɗanda suka ba da damar kafa \VUgras{shawa}
\textsuperscript{(7)} ba tare da tsoron cewa ruwan da yake tashi lokacin
amfani da ita zai lalata bangayen ba. \VUgras{Ƙaramin kwanon ƙarfe}
\textsuperscript{(8)}, wanda yake aikin wankan ƙafa, ya cika kayan
ɗakin.

\VUsaut{6.0mm}
\VUfilet{22mm}
